% P10-Core-Inhaltsprüfung, 15.09.2026: datierter Ergebnisanschluss; B-60 korrigiert.
% Belege: thesis-evidence/w2/core-content-audit-20260915/evaluation/
% P10-4b, 14.09.2026: Begriffs-/Beleganschlüsse und abgegrenzte Prüfaussagen präzisiert.
% P10-4a, 14.09.2026: Mechanismus, Quellenreichweite und v4-Nachweis angeglichen.
% M19 (12.09.2026): Leserführung und präzise Ergebnisdarstellung;
% Messwerte und historische Nachweisgrenzen bleiben erhalten.
% ============================================================
% Tabelle 7.1-1: Prüflandkarte — v01 (07.09.)
% In Overleaf ablegen als: tables/eval-pruefkarte.tex
% Aufruf im Fließtext: \input{\tabledir/eval-pruefkarte.tex}
% Label: t:eval-pruefkarte
% Stil: chap6.1-Konvention. Wird in 7.7 (t:eval-ergebnisprofil)
% mit Befunden geschlossen — P-Nummern dort identisch.
% HERKUNFT: anforderungsbezogene Neuformung der Claim-to-Evidence-
% Matrix (Schreibplan v2.4); Evidenzklassen [B]/[VA]/[D] =
% E2E-EVIDENZ.md §0.
% ============================================================

% M13: Prüflandkarte und Erläuterung dürfen eine Textseite teilen.
\begin{table}[!htb]
  \centering
  \small
  \begin{tabular}{l>{\raggedright\arraybackslash}p{0.47\textwidth}lcc}
    \hline
    & \textbf{Prüffrage} & \textbf{Anford.} & \textbf{Klasse} & \textbf{Ort} \\
    \hline
    P1 & Erfasst die Kette die fachlichen Aussagen der Quelle quellentreu und referenzgebunden? & A3 & [B] & \ref{sec:eval-kernzahlen} \\
    P2 & Welche unvollständig erfassten Referenzaussagen lassen sich den erzeugten Prüfhinweisen nachträglich zuordnen? & A3 & [B] & \ref{sec:eval-luecken} \\
    P3 & Erhält die interne Verdichtung (Kanonisierung) bereits erfasste Inhalte? & A4 & [B] & \ref{sec:eval-stufen} \\
    P4 & Werden im geprüften Fortschreibungspfad nur freigegebene Vorschläge übernommen, und blockiert die Speichernaht definierte ungültige Änderungen? & A4/A6 & [VA] & \ref{sec:eval-kontrollen} \\
    P5 & Werden neue, präzisierende, widersprechende, wiederholte und offene Informationen vertragsgemäß fortgeschrieben? & A5/A7 & [D] & \ref{sec:eval-fortschreibung} \\
    P6 & Welche Herkunftswege und Einzelreferenzen sind auflösbar, und welche aktuellen Änderungsbelege sind zugeordnet? & A3/A5 & [VA] & \ref{sec:eval-provenienz} \\
    P7 & Folgt die Außendarstellung (Issues/Dokumente) dem Core, und werden externe Eingänge kontrolliert zurückgeführt? & A8 & [D] & \ref{sec:eval-fortschreibung} \\
    P8 & Ist der Workflow über Prozessgrenzen fortsetzbar, und vermittelt die agentische Bedienschicht dieselben durch Gates kontrollierten Fähigkeiten? & A9/A10 & [D]+[VA] & \ref{sec:eval-resume} \\
    \hline
  \end{tabular}
  \caption[Prüffragen und zugeordnete Nachweise]{Prüflandkarte des Kapitels: acht Prüffragen zu den fünf
    anforderungsseitigen Eigenschaften mit Bezug auf die
    konsolidierten Designanforderungen A1--A10 (Kapitel~\ref{chap:explorative-problemklaerung}; A1/A2
    querschnittlich, s.\ Text), Evidenzklasse ([B] Benchmark gegen
    Referenzbestand, [VA] deterministischer Audit/Test,
    [D] dokumentierte Demonstration/Fallprüfung) und berichtendem Abschnitt.
    Abschnitt~\ref{sec:eval-ergebnisprofil} schließt die Landkarte mit Befunden und Grenzen.}
  \label{t:eval-pruefkarte}
\end{table}
